\section*{Capítulo \ref{ch:b3:complex-mechanisms} --- Mecanismos de reacción de los complejos}

\begin{solution}{exo:b3:complex-mechanisms:1}
Inertes: $\ce{[Cr(H2O)6]^{3+}}$ ($d^3$) y $\ce{[Co(NH3)6]^{3+}}$ ($d^6$ de bajo espín), con grandes energías de activación
del campo de ligandos. \omterm{def:b3:complex-mechanisms:labile}{Lábiles}: $\ce{[Cr(H2O)6]^{2+}}$ ($d^4$ de alto espín) y $\ce{[Cu(H2O)6]^{2+}}$ ($d^9$),
distorsionados por el efecto Jahn--Teller, con enlaces axiales largos y débiles; $\ce{[Zn(H2O)6]^{2+}}$ ($d^{10}$), sin barrera
de campo de ligandos.
\end{solution}

\begin{solution}{exo:b3:complex-mechanisms:2}
Cuando $k_2[\mathrm Y] \gg k_{-1}[\mathrm X]$, el denominador es $k_2[\mathrm Y]$ y $v = k_1[\mathrm{ML_5X}]$; cuando $k_{-1}[\mathrm X] \gg k_2[\mathrm Y]$, es
$k_{-1}[\mathrm X]$ y se obtiene la segunda forma (un preequilibrio inhibido por el ligando
saliente).
\end{solution}

\begin{solution}{exo:b3:complex-mechanisms:3}
D: $\Delta V^{\ddagger} > 0$ y $\Delta^{\ddagger}S^\circ > 0$ (un ligando se va). A: ambos negativos (un ligando está unido en el
estado de transición).
\end{solution}

\begin{solution}{exo:b3:complex-mechanisms:4}
$cis$-$\ce{[PtCl2(NH3)2]}$ a partir de $\ce{[PtCl4]^{2-}}$; $trans$-$\ce{[PtCl2(NH3)2]}$ a partir de $\ce{[Pt(NH3)4]^{2+}}$ (el
\omterm{def:b3:complex-mechanisms:trans-effect}{efecto trans} de \ce{Cl-} supera al de \ce{NH3}).
\end{solution}

\begin{solution}{exo:b3:complex-mechanisms:5}
$k_{\mathrm{obs}} = 2.0 \times \num{3.0e4} \times 0.10/1.2 = \qty{5.0e3}{s^{-1}}$; a \qty{1.0}{mol/L}, $\num{6.0e4}/3.0 = \qty{2.0e4}{s^{-1}}$; límite:
$k_i = \qty{3.0e4}{s^{-1}}$.
\end{solution}

\begin{solution}{exo:b3:complex-mechanisms:6}
$\Delta V^{\ddagger} = -RT\ln2/\Delta p = -8.314 \times 298 \times 0.693/\num{99.9e6} = \qty{-1.7e-5}{m^3.mol^{-1}} = \qty{-17}{cm^3.mol^{-1}}$:
asociativo.
\end{solution}

\begin{solution}{exo:b3:complex-mechanisms:7}
La fosfina, un $\sigma$-dador y un $\pi$-aceptor fuerte, debilita el enlace situado en trans respecto de ella (\omterm{def:b3:complex-mechanisms:trans-effect}{influencia
trans}, enlace más largo) y lo labiliza (\omterm{def:b3:complex-mechanisms:trans-effect}{efecto trans}): ese cloruro se
sustituye primero.
\end{solution}

\begin{solution}{exo:b3:complex-mechanisms:8}
El ligando puente acaba transferido al metal oxidado; la velocidad depende
mucho de la naturaleza del posible puente; se detecta un intermedio binuclear;
y la velocidad no puede superar la de la sustitución necesaria para formar el
puente sobre el compañero \omterm{def:b3:complex-mechanisms:labile}{lábil}.
\end{solution}

\begin{solution}{exo:b3:complex-mechanisms:9}
$(\lambda + \Delta G^\circ)^2/4\lambda$ con $4\lambda = 320$: 20.0, 11.3, 0 y \qty{5.0}{kJ/mol} (el último, en la \omterm{def:b3:complex-mechanisms:inverted}{región
invertida}).
\end{solution}

\begin{solution}{exo:b3:complex-mechanisms:10}
$k_{12} = \sqrt{4.0 \times \num{2.0e5} \times \num{1.0e4}} = \qty{8.9e4}{L.mol^{-1}.s^{-1}}$.
\end{solution}

\begin{solution}{exo:b3:complex-mechanisms:11}
$d^5$ de alto espín: 0; $d^7$: $1.33\,Dq$; $d^8$: $2.0\,Dq$. Velocidades de intercambio de agua esperadas: \ce{Mn^{2+}} $>$ \ce{Co^{2+}} $>$
\ce{Ni^{2+}}, el orden observado.
\end{solution}

\begin{solution}{exo:b3:complex-mechanisms:12}
En las reacciones muy exergónicas, la velocidad intrínseca supera la velocidad a la que los iones
se encuentran: la constante observada se queda en el límite de difusión y la disminución queda
oculta. Con el dador y el aceptor fijados a una distancia dada dentro de una molécula, no
hay etapa de difusión, y una serie de aceptores de fuerza impulsora creciente
mostró que la velocidad subía, pasaba por un máximo y bajaba.
\end{solution}

\begin{solution}{pb:b3:complex-mechanisms:1}
\textbf{1.} $5d^8$.
\textbf{2.} El campo fuerte de un ion $5d$ hace que el orbital $d_{x^2-y^2}$, que apunta hacia los cuatro
ligandos, quede muy alto: los ocho electrones se aparean en los cuatro orbitales inferiores.
\textbf{3.} El complejo tiene 16 electrones de valencia y una posición axial libre: el ligando
entrante se une primero, a través de un estado de transición pentacoordinado.
\textbf{4.} $k_{\mathrm{obs}} = k_1 + k_2[\mathrm Y]$: un camino a través del disolvente (ataque del disolvente y después sustitución rápida por Y)
y un camino directo.
\textbf{5.} La sustitución tarda horas (parte III): inerte a la escala de los minutos, no a la de
un día.
\textbf{6.} Cisplatino, porque el segundo \ce{NH3} sustituye a un cloruro situado en trans respecto de un cloruro
(\omterm{def:b3:complex-mechanisms:trans-effect}{efecto trans} $\ce{Cl-} > \ce{NH3}$); también se forman productos secundarios.
\textbf{7.} El yoduro tiene un \omterm{def:b3:complex-mechanisms:trans-effect}{efecto trans} mayor que el cloruro, lo que hace el efecto
director más limpio y más rápido.
\textbf{8.} En el $\ce{[PtI3(NH3)]-}$ se labiliza el yoduro situado en trans respecto de un yoduro, no el situado en trans respecto del
\ce{NH3}: el segundo amoníaco entra en cis respecto del primero.
\textbf{9.} Precipita yoduro de plata y queda en disolución $cis$-$\ce{[Pt(H2O)2(NH3)2]^{2+}}$.
\textbf{10.} El cloruro sustituye a las dos moléculas de agua en sus propias posiciones; los
ligandos ammina no se tocan.
\textbf{11.} A partir del $\ce{[Pt(NH3)4]^{2+}}$ y de dos cloruros.
\textbf{12.} Sus dos posiciones reactivas están opuestas: no puede unirse a dos
bases vecinas de la misma hebra de ADN, la lesión a través de la cual actúa el
cisplatino.
\textbf{13.} $\ln2/\num{9.0e-5} = \qty{7.7e3}{s} = \qty{2.1}{h}$.
\textbf{14.} $1 - \eu^{-\num{9.0e-5} \times 7200} = 0.48$.
\textbf{15.} El agua es un grupo saliente mucho mejor que el cloruro: el nitrógeno de una base del ADN
sustituye rápidamente al ligando acuo.
\textbf{16.} Asociativa, como todas las sustituciones del platino(II): \omterm{def:b3:complex-mechanisms:activation-volume}{volumen de activación}
negativo.
\textbf{17.} La alta concentración de cloruro desplaza el equilibrio de vuelta hacia la
forma dicloro, lo que mantiene intacto el fármaco hasta que llega a las células.
\textbf{18.} $f = K/(K + [\ce{Cl-}])$.
\textbf{19.} $\num{3.0e-3}/0.103 = 0.029$.
\textbf{20.} $\num{3.0e-3}/0.0070 = 0.43$.
\textbf{21.} 15.
\textbf{22.} Solo dentro de las células existe una fracción apreciable en forma de acuocomplejo
reactivo.
\textbf{23.} Los ligandos azufrados, como el glutatión y la cisteína y la metionina de las
proteínas, que se unen fuertemente al platino, blando.
\textbf{24.} El amoníaco se une fuertemente al platino y está en trans respecto del cloruro, un ligando de
\omterm{def:b3:complex-mechanisms:trans-effect}{efecto trans} pequeño: los enlaces con los ligandos ammina no se labilizan.
\textbf{25.} \textbf{La fracción del fármaco en la forma acuo reactiva es unas 15 veces mayor dentro de una
célula que en el plasma.}
\end{solution}
